% BHU PYQ -- Transcribed & Corrected
% Paper: COMMD21: Basic Accounting
% Exam: U.G. Semester-II Examination, 2024-25 (NEP)
% Department: Faculty of Commerce

\documentclass[11pt,a4paper]{article}
\usepackage[a4paper,top=2.5cm,bottom=2.5cm,left=2.5cm,right=2.5cm,headheight=14pt]{geometry}
\usepackage{amsmath,amssymb}
\usepackage[T1]{fontenc}
\usepackage[utf8]{inputenc}
\usepackage{lmodern,microtype,fancyhdr,enumitem,hyperref,lastpage,parskip,booktabs}

\hypersetup{colorlinks=true,urlcolor=black,linkcolor=black,pdftitle={COMMD21: Basic Accounting -- BHU Sem II 2024-25 (NEP MDC)}}
\pagestyle{fancy}\fancyhf{}
\fancyhead[L]{\small   \textit{Banaras Hindu University}}
\fancyhead[R]{\small   \textit{COMMD21: Basic Accounting}}
\fancyfoot[C]{\small Page  \thepage\ of \pageref{LastPage}}


\newlist{parts}{enumerate}{2}
\setlist[parts,1]{label=(\alph*),leftmargin=*,itemsep=4pt,topsep=2pt}
\setlist[parts,2]{label=(\roman*),leftmargin=*,itemsep=3pt,topsep=2pt}


\begin{document}


\begin{center}
  {\large\bfseries BANARAS HINDU UNIVERSITY}\\[2pt]
  {
\normalsize U.G. Semester-II Examination, 2024-25 (NEP)}\\[6pt]
  \rule{0.8\linewidth}{0.5pt}\\[6pt]
  {\large\bfseries FACULTY OF COMMERCE}\\[4pt]
  {\large\bfseries Paper Code: COMMD21 --- Basic Accounting}\\[4pt]
  {
\normalsize Time: 2:30 Hours\hfill Maximum Marks: 70}
\end{center}

\rule{\linewidth}{0.4pt}\smallskip



\noindent   \textit{Instruction: Write your Roll No. at the top immediately on the receipt of this question paper.}\\[2pt]



\noindent   \textit{Note: Answer all questions. All questions carry equal marks.}
\bigskip




\noindent   \textbf{Question 1}

\begin{parts}
  \item Differentiate between Book Keeping and Accountancy.
  \item Explain in brief various Concepts and Conventions of Accounting.
\end{parts}


\begin{center}   \textbf{OR}\end{center}

Define Accounting. Explain the characteristics, functions, process and limitations of Accounting.

\medskip\rule{\linewidth}{0.2pt}\medskip




\noindent   \textbf{Question 2} \\
Record the following transactions in Journal:

\begin{center}
\small

\begin{tabular}{r p{8cm} r}
 oprule
   \textbf{Date (Jan 2025)} &   \textbf{Particulars} &   \textbf{Amount (Rs.)} \\
\midrule
Jan 1 & Kumar commenced business with cash & 7,000 \\
Jan 2 & Purchased furniture & 70 \\
Jan 3 & Purchased goods from Prakash & 4,000 \\
Jan 4 & Goods sold to Ram Lal & 5,000 \\
Jan 5 & Goods sold to Prem & 6,500 \\
Jan 6 & Goods purchased from Ramesh & 5,000 \\
Jan 7 & Goods sold to Ashok & 8,000 \\
Jan 8 & Ashok returned goods & 500 \\
Jan 10 & Paid Wages & 100 \\
Jan 14 & Received from Ram Lal in full settlement & 4,950 \\
Jan 15 & Goods sold to Amar for cash & 4,000 \\
Jan 20 & Withdrew for personal use & 1,000 \\
Jan 22 & Paid to Prakash & 20,000 \\
Jan 25 & Received from Prem & 5,000 \\
Jan 26 & Cash Sales & 8,000 \\
Jan 29 & Paid rent for month & 2,000 \\
Jan 31 & Paid salary & 3,000 \\
ottomrule
\end{tabular}
\end{center}


\begin{center}   \textbf{OR}\end{center}

From the following transactions, prepare a Triple Column Cash Book:

\begin{center}
\small

\begin{tabular}{r p{8.5cm} r}
 oprule
   \textbf{Date (July 2024)} &   \textbf{Particulars} &   \textbf{Amount (Rs.)} \\
\midrule
July 1 & Commenced business with cash & 50,000 \\
July 2 & Opened Bank Account & 20,000 \\
July 3 & Bought Furniture & 4,000 \\
July 4 & Purchased goods & 25,000 \\
July 5 & Sold goods & 12,000 \\
July 9 & Purchased goods by Cheque & 5,000 \\
July 10 & Cash sales deposited into Bank & 3,000 \\
July 15 & Withdrew cash for business & 2,000 \\
July 17 & Paid electricity bill by Cheque & 1,500 \\
July 18 & Received commission by Cheque & 500 \\
July 20 & Bought goods from Rohan & 7,000 \\
July 24 & Sold goods to Mohan & 9,000 \\
July 26 & Received a cheque from Mohan in full settlement & 8,800 \\
July 27 & Paid to Sohan in full settlement & 6,850 \\
July 28 & Deposited cash into Bank & 1,200 \\
July 31 & Paid rent by cheque Rs. 1,300 and salary in cash & 1,500 \\
ottomrule
\end{tabular}
\end{center}

\medskip\rule{\linewidth}{0.2pt}\medskip




\noindent   \textbf{Question 3} \\
What are the various types of errors in Accounts? Explain with imaginary examples.


\begin{center}   \textbf{OR}\end{center}

Draw a Trial Balance from the following balances extracted from ledger accounts of M/s Brothers as on 31 Dec., 2024:

\begin{center}
\small

\begin{tabular}{p{4.5cm} r p{4.5cm} r}
 oprule
   \textbf{Account} &   \textbf{Amount (Rs.)} &   \textbf{Account} &   \textbf{Amount (Rs.)} \\
\midrule
Capital & 15,000 & Advertisement & 210 \\
Land and Building & 15,600 & Rent and Taxes & 160 \\
Bank Overdraft & 2,500 & Insurance & 40 \\
Cash in Hand & 680 & Discount Allowed & 300 \\
Opening Stock & 6,000 & Repairs & 210 \\
Purchases & 7,200 & Interest Received & 500 \\
Sales & 17,000 & Debtors & 6,620 \\
Provision for Bad Debts & 370 & Creditors & 4,100 \\
Wages & 1,250 & General Expenses & 500 \\
Salary & 700 & & \\
ottomrule
\end{tabular}
\end{center}

\medskip\rule{\linewidth}{0.2pt}\medskip




\noindent   \textbf{Question 4}

\begin{parts}
  \item Differentiate between Trial Balance and Balance Sheet.
  \item Show the proforma of Trading Account, Profit and Loss Account and Balance Sheet separately.
\end{parts}


\begin{center}   \textbf{OR}\end{center}

Prepare Trading and Profit and Loss Account of M/s Naresh for the year ending 31 Dec., 2024 and Balance Sheet on that date from the following Trial Balance:

\begin{center}
\small

\begin{tabular}{p{4.5cm} r p{4.5cm} r}
 oprule
   \textbf{Debit Particulars} &   \textbf{Amount (Rs.)} &   \textbf{Credit Particulars} &   \textbf{Amount (Rs.)} \\
\midrule
Cash in Hand & 10,000 & Bank Overdraft & 10,000 \\
Opening Stock & 25,000 & Capital & 80,000 \\
Purchases & 75,000 & Interest & 2,000 \\
Sales Return & 5,000 & Sales & 1,61,000 \\
Salary & 15,000 & Purchase Return & 2,500 \\
Debtors & 1,00,000 & Commission & 2,500 \\
Machinery & 30,000 & Creditors & 60,000 \\
Wages & 10,000 & Bills Payable & 15,700 \\
Bad Debts & 3,500 & Discount & 800 \\
Furniture & 10,000 & General Reserve & 12,000 \\
Carriage Inwards & 2,000 & & \\
Office Expenses & 6,000 & & \\
Insurance & 6,000 & & \\
Bills Receivables & 3,500 & & \\
Printing and Stationery & 4,000 & & \\
Rent & 7,500 & & \\
Drawings & 12,000 & & \\
Land & 22,000 & & \\
\midrule
   \textbf{Total} &   \textbf{3,46,500} &   \textbf{Total} &   \textbf{3,46,500} \\
ottomrule
\end{tabular}
\end{center}




\noindent   \textbf{Adjustments:}

\begin{enumerate}[leftmargin=*,itemsep=2pt,topsep=3pt]
  \item Outstanding salary = Rs. 1,500 and outstanding wages = Rs. 1,200.
  \item Prepaid Insurance = Rs. 1,000.
  \item Accrued Interest = Rs. 250.
  \item Depreciate Furniture @ 10\% p.a. and Machinery @ 15\% p.a.
  \item Unexpired commission = Rs. 500.
  \item Closing Stock = Rs. 35,000.
\end{enumerate}

\medskip\rule{\linewidth}{0.2pt}\medskip




\noindent   \textbf{Question 5} \\
Write short notes on the following:

\begin{parts}
  \item Users of Accounting
  \item Preparation of Journal
  \item Rectification of Errors
  \item Adjustments
\end{parts}

\vfill\rule{\linewidth}{0.4pt}

\begin{center}\small
\href{https://bhu-exam.vercel.app/}{Click here to visit our website}\\[4pt]
\href{https://me-aryan.vercel.app/}{Click here to visit our contributor}
\end{center}

\end{document}
